Small instruction-tuned language models sometimes answer ``unknown'' even when their context contains enough
evidence to answer correctly. In Qwen2.5-3B-Instruct, this false abstention can follow a note in which the model
has already quoted the relevant fact. Earlier interventions recovered correct answers but also encouraged
unsupported answers. We ask whether an agent-proposed weight edit can improve this trade-off.
Following weight-space programming \citep{takehi2026wsp}, a Claude agent inspected the model on an exploration
split and proposed three edits. We evaluated them in a pre-registered case study in one model, with separate
selection and sealed confirmation splits, a transfer split, and six frozen comparators serving repair or pathway
tests. The selected repair was a rank-1 change to the attention output projection in layer 35. It met the
pre-registered repair criterion on sealed data: the base-weighted recovery rate on false abstentions was 82\%,
while unsupported answers on unanswerable questions rose from 3.5\% to 4.2\%. It damaged already-correct answers
at a base-weighted rate of 3\%. Its release-versus-unsupported-answer trade-off was better than a prompt
instruction, activation steering and directional orthogonalisation. The comparison with a simple logit bias
against ``unknown'' was inconclusive. A separate two-head edit provided evidence about a pathway through which
answers can be released, but the agent's explanation of that component was not supported by the pre-registered
tests.
